% --- Avant-propos -----------------------------------------------------------
\chapter*{Avant-propos}
\markboth{Avant-propos}{Avant-propos}
\addcontentsline{toc}{chapter}{Avant-propos}

Ce livre est né d'un constat simple : entre le programme officiel, qui dit
\emph{ce qu'il faut} enseigner, et la salle de classe, où il faut décider
\emph{comment} le faire tenir en une heure devant quarante élèves, il manque
un document. Les répartitions annuelles donnent le calendrier, les documents
d'accompagnement donnent l'esprit, mais aucun des deux ne met entre les mains
du professeur et de l'élève le même objet, ouvert à la même page.

C'est ce qu'on a essayé de faire ici. Chaque chapitre est écrit deux fois dans
un seul texte : l'élève y trouve un cours qu'il peut lire seul le soir, le
professeur y trouve les corrigés, les erreurs qu'il verra passer et les
remarques de conduite de séance. Les deux éditions sortent du même fichier
source ; rien ne peut donc diverger entre ce que l'un explique et ce que
l'autre révise.

\section*{Un mot à l'élève de première}

L'année de première est celle où les objets cessent d'être donnés et
commencent à être construits. En seconde, une fonction était une formule dont
on calculait des images ; ici, elle a un ensemble de définition qu'il faut
établir, des limites aux bornes de cet ensemble, une dérivée dont le signe
gouverne les variations, et une courbe que l'on trace parce qu'on a démontré
son allure, non parce qu'on a relié des points. La même bascule se produit
partout : le vecteur devient un outil de calcul, l'expérience aléatoire un
ensemble que l'on dénombre, la proposition une chose dont on discute la
vérité.

C'est aussi l'année où l'on prend de l'avance. Presque tout ce qui vous
attend en terminale s'appuie sur les pages qui suivent : la dérivation, les
suites, le barycentre, le dénombrement. Ce que vous installerez solidement
cette année, vous ne le réapprendrez pas l'an prochain — vous vous en
servirez.

Vous remarquerez que chaque chapitre commence par une page de révision qui
ne parle jamais de la notion nouvelle. C'est voulu. On ne vous demandera
jamais de deviner ce que vous n'avez pas encore appris : on vérifie d'abord
que les outils de l'an dernier répondent encore, puis on ouvre le cours.
Prenez ces pages au sérieux, elles coûtent dix minutes et vous en font gagner
beaucoup.

\section*{Un mot au professeur}

Le Programme d'Études attribue à la classe de T11 cinq heures de
mathématiques par semaine, soit cent cinquante heures réparties sur quatre
composantes : l'analyse, l'algèbre, la géométrie et le traitement des données
et probabilités. Le livre suit ces quatre composantes ; la table de
progression annuelle, en annexe, indique comment les vingt chapitres se
placent dans les cinq périodes du calendrier scolaire et ce qui peut être
allégé lorsqu'une période est amputée.

Il n'y a pas d'examen national au terme de cette année, et le livre ne fait
donc répéter aucune épreuve. Les exercices repris d'examens viennent des
sujets de fin de première d'autres systèmes francophones, qui interrogent les
mêmes notions ; chacun porte sa provenance. Le bloc qui ferme chaque chapitre
va chercher, parmi les questions de baccalauréat, celles qui ne réclament que
ce qui vient d'être enseigné : c'est le moyen le plus honnête de montrer à
une classe de première ce qui l'attend, sans lui demander ce qu'elle ne peut
pas encore faire.

Les démonstrations sont toutes écrites. Cela ne veut pas dire qu'elles doivent
toutes être faites au tableau : à vous de choisir. Elles sont là pour que
l'élève qui veut comprendre pourquoi un théorème est vrai trouve la réponse
dans son livre plutôt que nulle part.

\vspace{1.2em}
\noindent\textit{Antananarivo}
